\chapter{An\'alisis y requisitos}\label{cap_requisitos}
% tfg-floatbarrier-auto
\makeatletter
\@ifundefined{tfgfloatbarrier}{
  \def\tfgfloatbarrier{1}
  \let\tfgorigfigure\figure
  \let\tfgendorigfigure\endfigure
  \renewenvironment{figure}{\tfgorigfigure}{\tfgendorigfigure\FloatBarrier}
  \let\tfgorigtable\table
  \let\tfgendorigtable\endtable
  \renewenvironment{table}{\tfgorigtable}{\tfgendorigtable\FloatBarrier}
}{}
\makeatother

\FloatBarrier

\section{El proceso de an\'alisis en la Ingenier\'ia del Software}

El an\'alisis de requisitos es una etapa fundamental dentro del ciclo de vida de la Ingenier\'ia del Software. Su prop\'osito es definir con precisi\'on las necesidades del dominio del problema antes de plantear una soluci\'on tecnol\'ogica. Un an\'alisis exhaustivo previene sobrecostes, minimiza el riesgo de desarrollo de funcionalidades innecesarias y asegura que el sistema entregado se alinee con las expectativas de los usuarios.

En el contexto de este proyecto, la metodolog\'ia de an\'alisis ha partido de un estudio minucioso de la situaci\'on actual en el \'ambito docente, identificando las limitaciones operativas existentes. 

\subsection{Situaci\'on actual y necesidades de negocio}
El escenario de partida presenta una fricci\'on operativa notable: creaci\'on de actividades y recogida de entregas con alto componente manual, dificultad para mantener trazabilidad uniforme y baja capacidad de evoluci\'on. Las necesidades de negocio priorizadas son reducir la carga administrativa, ofrecer al estudiante un flujo de entrega m\'as guiado, centralizar evaluaci\'on y mejorar la mantenibilidad del sistema. Estas necesidades conectan directamente con los bloques funcionales del MVP.

\FloatBarrier
\section{Especificaci\'on de Requisitos del Sistema (ERS)}

\subsection{Fundamentaci\'on te\'orica del ERS}
La Especificaci\'on de Requisitos del Sistema (ERS) es el documento est\'andar que act\'ua como contrato entre las necesidades del negocio y la implementaci\'on t\'ecnica. De forma gen\'erica, su objetivo es catalogar qu\'e debe hacer el sistema (requisitos funcionales) y bajo qu\'e restricciones de calidad y seguridad debe operar (requisitos no funcionales y reglas de negocio). 

Para este Trabajo Fin de Grado en particular, el ERS tiene como objetivo espec\'ifico delimitar de forma inquebrantable qu\'e operaciones est\'an permitidas para los roles de profesor, estudiante y administrador, as\'i como definir los atributos esenciales de los cursos, actividades y entregables.



\subsection{Fuentes de requisitos}

La base fundamental para la definici\'on de los requisitos del proyecto es el propio Documento ERS (Especificaci\'on de Requisitos del Sistema), el cual constituye de por s\'i el cat\'alogo formal y exhaustivo de requisitos. Este documento fue redactado, modelado estructuralmente y versionado haciendo uso de la herramienta Proteus \cite{proteus}, lo que garantiza rigor metodol\'ogico en su definici\'on.

A partir de este cat\'alogo central establecido en el ERS, el flujo de desarrollo se ha complementado con el Documento DAS (Documento de An\'alisis del Sistema) para el modelado conceptual, as\'i como con el uso din\'amico de un Product Backlog y un Sprint Backlog durante la ejecuci\'on iterativa. La interacci\'on continua con estas herramientas, junto con la evidencia emp\'irica obtenida en las pruebas funcionales y de integraci\'on, asegura que los requisitos no se mantengan como meras definiciones abstractas. En lugar de una especificaci\'on est\'atica, los requisitos del ERS se priorizan, se integran en incrementos de valor y se validan con el comportamiento del producto en funcionamiento, evitando en todo momento cualquier desconexi\'on entre la especificaci\'on te\'orica y la implementaci\'on real del sistema.

\subsection{Estructura del cat\'alogo}

El cat\'alogo se organiza en cuatro bloques complementarios: requisitos funcionales de informaci\'on y operaci\'on, reglas de negocio que acotan el comportamiento permitido, requisitos no funcionales de calidad del sistema y restricciones e integraciones t\'ecnicas por entorno. Esta estructura permite cubrir tanto qu\'e debe hacer el sistema como bajo qu\'e condiciones debe hacerlo.

En Proteus \cite{proteus} se distinguen expl\'icitamente requisitos funcionales (por ejemplo,
\texttt{RQF-001}), reglas de negocio (\texttt{REGN-001}) y no funcionales
(\texttt{RQNF-001} a \texttt{RQNF-011}), lo que facilita su seguimiento
durante la implementaci\'on.

\subsection{Requisitos funcionales}

El ERS define un abanico funcional amplio, del que se ha priorizado el n\'ucleo
necesario para el MVP. A nivel de dominio, la funcionalidad se agrupa en identidad y acceso (autenticaci\'on de usuarios y control por rol), estructura acad\'emica (cursos, grupos, actividades y vinculaciones), gesti\'on de entregables (configuraci\'on de fechas, visibilidad, reenv\'ios y estructura esperada), gesti\'on de entregas (subida, historial, descarga y trazabilidad) y evaluaci\'on y feedback (calificaci\'on, comentarios y publicaci\'on de notas). Esta agrupaci\'on garantiza cobertura del ciclo acad\'emico completo.

Como ejemplos representativos se incluyen \textbf{RQF-001} (creaci\'on de usuarios por administrador), \textbf{RQF-014--RQF-017} (operaciones de entrega y versionado) y \textbf{RQF-018--RQF-025} (evaluaci\'on, visibilidad e integraciones). Estos ejemplos reflejan la amplitud funcional sin agotar el inventario completo.

\subsection{Reglas de negocio}

Las reglas de negocio restringen de forma expl\'icita qu\'e acciones son admisibles.
Una regla clave es \textbf{REGN-001}, que establece que solo el profesor
autorizado puede asignar una calificaci\'on.

Esta norma se implementa en validaciones de servicio y controlador,
evitando tanto la elevaci\'on de privilegios como la evaluaci\'on fuera de contexto
(p. ej., docente no vinculado al grupo o actividad).

\subsection{Requisitos no funcionales}

Los requisitos no funcionales aportan condiciones de calidad del producto.
La Tabla~\ref{tab:rqnf-resumen} resume los m\'as relevantes para el alcance actual.

\begin{table}[!htbp]
\centering
\begin{tabular}{|P{2.1cm}|P{3.5cm}|P{6.4cm}|}
\hline
\textbf{C\'odigo} & \textbf{Categor\'ia} & \textbf{Aplicaci\'on en el sistema} \\
\hline
RQNF-001 & Fiabilidad & Mitigar p\'erdida de datos en operaciones cr\'iticas de entrega/evaluaci\'on. \\
\hline
RQNF-002 & Usabilidad & Flujos entendibles para usuarios no t\'ecnicos, con navegaci\'on guiada por rol. \\
\hline
RQNF-004 & Mantenibilidad & Bajo acoplamiento entre capas y separaci\'on estricta de responsabilidades. \\
\hline
RQNF-006 & Eficiencia & Respuesta \`agil en consultas frecuentes de panel, actividad y entregables. \\
\hline
RQNF-008 & Portabilidad & Compatibilidad web en navegadores modernos de uso general. \\
\hline
RQNF-009 & Seguridad & Exposici\'on de datos condicionada por autenticaci\'on y autorizaci\'on. \\
\hline
\end{tabular}
\caption{Requisitos no funcionales representativos}
\label{tab:rqnf-resumen}
\end{table}

\subsection{Priorizaci\'on y alcance}

La priorizaci\'on se realiz\'o con criterio MoSCoW \cite{moscow_method} y validaci\'on continua con el backlog.
El foco del MVP se situ\'o en los procesos con mayor impacto docente: gesti\'on de actividades y entregables, flujo completo de entrega y evaluaci\'on, seguridad por rol, y persistencia con trazabilidad de cambios. Esta selecci\'on maximiza valor operativo en el horizonte temporal disponible.

Funcionalidades avanzadas (integraciones extendidas, autenticaci\'on federada completa,
observabilidad m\'as profunda) se mantienen en backlog de evoluci\'on y no comprometen
la viabilidad del alcance de TFG.

\subsection{Trazabilidad requisito-backlog-c\'odigo}

La trazabilidad se gestiona en tres niveles estrechamente interconectados: en primer lugar, un \textbf{nivel de an\'alisis} apoyado por el identificador un\'ivoco establecido en el ERS a trav\'es de Proteus; en segundo lugar, un \textbf{nivel de planificaci\'on} guiado por los PBI (\emph{Product Backlog Items}, elementos de trabajo discretos que representan funcionalidades con valor para el usuario dentro de las metodolog\'ias \`agiles) y su sprint asociado; y por \'ultimo, un \textbf{nivel de ejecuci\'on} que liga estos elementos al m\'odulo de software implementado y a la prueba automatizada que lo valida. Esta triple vista permite justificar cualquier decisi\'on de dise\~no con evidencia t\'ecnica medible.

La Tabla~\ref{tab:trazabilidad-requisitos} muestra ejemplos representativos de esta trazabilidad, ilustrando c\'omo un requisito formal del ERS se convierte en un PBI dentro de un sprint de desarrollo y, finalmente, se materializa en c\'odigo y pruebas espec\'ificas.

\begin{table}[!htbp]
\centering
\begin{tabular}{|P{2.2cm}|P{3.2cm}|P{4.0cm}|P{2.6cm}|}
\hline
\textbf{Req.} & \textbf{PBI / Sprint} & \textbf{Implementaci\'on principal} & \textbf{Validaci\'on} \\
\hline
RQF-001 & Usuarios y roles / S1-S2 & AuthController, UsuarioService, seguridad por rol & Tests auth + integraci\'on \\
\hline
RQF-018 & Crear entregable / S3-S4 & EntregableController, EntregableService, Entregable & Unit + API tests \\
\hline
RQF-023 & Subir entrega / S4-S5 & EntregaController, EntregaService, MaterialService & Integraci\'on + E2E \\
\hline
REGN-001 & Evaluaci\'on docente / S5 & Validaciones en servicio de entregas y evaluaci\'on & Pruebas de autorizaci\'on \\
\hline
RQNF-009 & Seguridad global / transversal & SecurityConfig, JwtAuthenticationFilter, RateLimitFilter & Tests de seguridad \\
\hline
\end{tabular}
\caption{Ejemplo de trazabilidad entre requisitos, backlog e implementaci\'on}
\label{tab:trazabilidad-requisitos}
\end{table}

\subsection{Contexto documental de partida (ERS y DAS)}

La especificaci\'on documental de este TFG no se apoya en descripciones ad hoc, sino en dos artefactos de an\'alisis formal mantenidos en Proteus. El \textbf{ERS (Especificaci\'on de Requisitos del Sistema)} define situaci\'on actual, necesidades de negocio, subsistemas y cat\'alogo de requisitos (generales, de informaci\'on, funcionales, no funcionales y reglas de negocio), mientras que el \textbf{DAS (Documento de An\'alisis del Sistema)} define arquitectura l\'ogica, modelo de clases, diagramas de estado, diagramas de secuencia, navegabilidad e inventario de mockups de interfaz. Esta base asegura continuidad documental entre an\'alisis y ejecuci\'on.

La reutilizaci\'on de ambos documentos en esta memoria persigue dos objetivos:
primero, preservar trazabilidad respecto al trabajo original de Ingenier\'ia
de Requisitos; segundo, demostrar que las decisiones de dise\~no e implementaci\'on
adoptadas durante el TFG responden a una base de especificaci\'on expl\'icita.



El ERS parte de un an\'alisis detallado de un escenario acad\'emico real donde los procesos docentes presentaban una alta ineficiencia: los profesores requer\'ian m\'ultiples pasos manuales y plataformas dispersas para la definici\'on de actividades, publicaci\'on de material y, de manera muy acusada, para la recogida, versionado y almacenamiento de las entregas de los alumnos. El rastreo de estas entregas carec\'ia de una trazabilidad uniforme (fechas, reenv\'ios, visibilidad), lo que derivaba en p\'erdidas de informaci\'on y en un soporte t\'ecnico que imposibilitaba la integraci\'on de l\'ogicas de evaluaci\'on m\'as avanzadas.

Ante este diagn\'ostico espec\'ifico, las necesidades de negocio que fundamentan este TFG se orientan a un cambio de paradigma centralizado. El objetivo es proporcionar un entorno \'unico donde la carga administrativa del profesor se reduzca mediante flujos estandarizados (desde la creaci\'on de un \emph{entregable} hasta la introducci\'on de calificaciones), y donde el alumno perciba transparencia absoluta sobre sus fechas l\'imite, el estado de sus versiones y los comentarios (feedback) recibidos. Adem\'as, era imperativo garantizar que la evaluaci\'on solo pudiera ser efectuada bajo reglas de autorizaci\'on y contexto acad\'emico estrictas (vinculaci\'on estudiante-grupo-profesor). 

Estas necesidades condicionan el resto de los cap\'itulos en tres planos tangibles: la definici\'on arquitect\'onica (separando estrictamente los servicios de usuarios de la gesti\'on acad\'emica), la implementaci\'on a trav\'es de m\'odulos enfocados a dominios espec\'ificos, y una bater\'ia de pruebas concebida para validar tanto los flujos de experiencia de usuario como el cumplimiento inquebrantable de las reglas de autorizaci\'on impuestas en el negocio.

\subsection{Descomposici\'on en subsistemas (ERS)}

El ERS descompone el sistema en cinco grandes subsistemas l\'ogicos para organizar sus funcionalidades. La Tabla \ref{tab:subsistemas-ers} detalla esta descomposici\'on, la cual ha sido reutilizada como gu\'ia fundamental de dise\~no modular a lo largo del desarrollo del TFG.

\begin{table}[!htbp]
\centering
\small
\begin{tabular}{|P{2.0cm}|P{3.9cm}|P{5.9cm}|}
\hline
\textbf{C\'odigo} & \textbf{Subsistema} & \textbf{Responsabilidad principal} \\
\hline
SUBS-001 & SIU (Interacci\'on con el Usuario) & Portal web y vistas por rol para alumno, profesor y administrador. \\
\hline
SUBS-002 & SGUS (Usuarios y Seguridad) & Autenticaci\'on, autorizaci\'on por rol y control de acceso contextual. \\
\hline
SUBS-003 & SGA (Gesti\'on Acad\'emica) & Cursos, actividades, entregables, evaluaci\'on y feedback. \\
\hline
SUBS-004 & SPA (Procesamiento de Archivos) & Validaci\'on t\'ecnica de ficheros, organizaci\'on y versionado de entregas. \\
\hline
SUBS-005 & SPD (Persistencia de Datos) & Persistencia relacional y repositorio de archivos con trazabilidad. \\
\hline
\end{tabular}
\normalsize
\caption{Subsistemas definidos en ERS y reutilizados en el dise\~no del TFG}
\label{tab:subsistemas-ers}
\end{table}

\subsection{Cat\'alogo exhaustivo de requisitos del ERS}

Con objeto de incorporar de forma completa el contenido de ERS, se documenta
el inventario de requisitos por tipolog\'ia y su interpretaci\'on en el alcance
real del TFG.

\FloatBarrier
\subsubsection*{Requisitos generales (RQG)}

La Tabla \ref{tab:rqg-inventario} detalla los requisitos generales del sistema, que establecen el marco de operaci\'on b\'asico para la gesti\'on de usuarios, actividades, entregables y la comunicaci\'on.

\begin{table}[!htbp]
\centering
\small
\begin{tabular}{|P{1.4cm}|P{3.2cm}|P{7.2cm}|}
\hline
\textbf{C\'odigo} & \textbf{Nombre} & \textbf{Descripci\'on sint\'etica} \\
\hline
RQG-001 & Gesti\'on de usuarios y roles & Alta de usuarios con rol (administrador, profesor, alumno) y permisos diferenciados por perfil. \\
\hline
RQG-002 & Gesti\'on de actividades y entregables & Modelado de actividades evaluables/no evaluables y entregables con reglas propias. \\
\hline
RQG-003 & Acceso de alumnos a entregables & Acceso de estudiante a entregables visibles y registro temporal de acciones. \\
\hline
RQG-004 & Gesti\'on de versiones y reenv\'ios & Control de versiones de entrega y soporte de reenv\'io cuando el docente lo habilita. \\
\hline
RQG-005 & Configuraci\'on de actividades y entregables & Definici\'on de fechas, formatos y condiciones de publicaci\'on por contexto acad\'emico. \\
\hline
RQG-006 & Comunicaci\'on y feedback & Canal de comunicaci\'on docente-estudiante vinculado a la evaluaci\'on de entregas. \\
\hline
RQG-007 & Seguridad y control de acceso & Restricci\'on de visibilidad/operaci\'on por identidad, rol y pertenencia a curso-grupo. \\
\hline
\end{tabular}
\normalsize
\caption{Inventario de requisitos generales (RQG) en ERS}
\label{tab:rqg-inventario}
\end{table}

\FloatBarrier
\subsubsection*{Requisitos de informaci\'on (RQI)}

La Tabla \ref{tab:rqi-inventario} describe los requisitos de informaci\'on que definen qu\'e datos esenciales deben persistirse en el sistema (actividades, entregables, usuarios, etc.).

\begin{table}[!htbp]
\centering
\small
\begin{tabular}{|P{1.4cm}|P{3.2cm}|P{7.2cm}|}
\hline
\textbf{C\'odigo} & \textbf{Nombre} & \textbf{Descripci\'on sint\'etica} \\
\hline
RQI-001 & Informaci\'on sobre actividades & Persistencia de metadatos de actividades creadas por profesorado. \\
\hline
RQI-002 & Informaci\'on sobre entregables & Persistencia de configuraci\'on de cada entregable y su vinculaci\'on con actividad. \\
\hline
RQI-003 & Informaci\'on de usuarios & Registro de identidad, rol y datos de operaci\'on por tipo de usuario. \\
\hline
RQI-004 & Informaci\'on de cursos & Persistencia de estructura acad\'emica (curso/grupo y relaciones). \\
\hline
RQI-005 & Informaci\'on de feedback & Persistencia de calificaciones, comentarios y estado de publicaci\'on. \\
\hline
\end{tabular}
\normalsize
\caption{Inventario de requisitos de informaci\'on (RQI) en ERS}
\label{tab:rqi-inventario}
\end{table}

\subsubsection*{Requisitos funcionales (RQF)}

La Tabla \ref{tab:rqf-inventario-1} y la Tabla \ref{tab:rqf-inventario-2} recogen, de manera estructurada y en dos bloques para facilitar su lectura, las capacidades operativas que el sistema debe ofrecer al usuario final, desde la administraci\'on de identidades hasta la integraci\'on cloud.

\begin{table}[!htbp]
\centering
\small
\begin{tabular}{|P{1.4cm}|P{3.3cm}|P{7.1cm}|}
\hline
\textbf{C\'odigo} & \textbf{Nombre} & \textbf{Descripci\'on sint\'etica} \\
\hline
RQF-001 & Crear usuarios & Alta de usuarios por administrador con asignaci\'on de rol. \\
\hline
RQF-002 & Iniciar sesi\'on & Autenticaci\'on de usuarios en la plataforma. \\
\hline
RQF-003 & Editar usuarios & Modificaci\'on de datos de usuario seg\'un permisos. \\
\hline
RQF-004 & Eliminar usuarios & Baja de usuarios por perfil administrador. \\
\hline
RQF-005 & Crear actividades & Creaci\'on de actividad acad\'emica por docente. \\
\hline
RQF-006 & Visibilizar actividades & Publicaci\'on/ocultaci\'on de actividades para estudiantes. \\
\hline
RQF-007 & Proporcionar adjuntos & Asociaci\'on de material de apoyo a actividad/entregable. \\
\hline
RQF-008 & Modificar actividades & Edici\'on de configuraci\'on de actividad por docente propietario. \\
\hline
RQF-009 & Eliminar actividades & Eliminaci\'on de actividad seg\'un reglas de propiedad. \\
\hline
RQF-010 & Ver actividades & Consulta de actividades por rol y contexto acad\'emico. \\
\hline
RQF-011 & Crear entregables & Alta de entregables dentro de actividad. \\
\hline
RQF-012 & Modificar entregables & Edici\'on de entregable por docente autorizado. \\
\hline
\end{tabular}
\normalsize
\caption{Inventario de requisitos funcionales (RQF-001 a RQF-012)}
\label{tab:rqf-inventario-1}
\end{table}

\begin{table}[!htbp]
\centering
\small
\begin{tabular}{|P{1.4cm}|P{3.3cm}|P{7.1cm}|}
\hline
\textbf{C\'odigo} & \textbf{Nombre} & \textbf{Descripci\'on sint\'etica} \\
\hline
RQF-013 & Eliminar entregable & Baja de entregable por docente autorizado. \\
\hline
RQF-014 & Realizar entrega & Subida de archivo por estudiante para un entregable habilitado. \\
\hline
RQF-015 & Enviar actividad & Confirmaci\'on de env\'io de actividad completa. \\
\hline
RQF-016 & Versionar entregables & Mantenimiento de historial de versiones por entregable. \\
\hline
RQF-017 & Versionar actividades & Mantenimiento de versiones de resoluci\'on de actividad. \\
\hline
RQF-018 & Dar feedback & Publicaci\'on de comentarios docentes sobre entrega/actividad. \\
\hline
RQF-019 & Realizar calificaci\'on & Asignaci\'on de nota por profesor autorizado. \\
\hline
RQF-020 & Ver entregables & Consulta de entregas y versiones por rol permitido. \\
\hline
RQF-021 & Crear curso & Alta de curso/asignatura por administraci\'on. \\
\hline
RQF-022 & Modificar curso & Edici\'on de datos de curso por administraci\'on. \\
\hline
RQF-023 & Vista de estudiante & Visualizaci\'on por profesor de la experiencia de alumno. \\
\hline
RQF-024 & Visibilizar entregable & Publicaci\'on/ocultaci\'on de entregables concretos. \\
\hline
RQF-025 & Integraci\'on cloud & Uso de almacenamiento en nube seg\'un configuraci\'on de entorno. \\
\hline
\end{tabular}
\normalsize
\caption{Inventario de requisitos funcionales (RQF-013 a RQF-025)}
\label{tab:rqf-inventario-2}
\end{table}

\FloatBarrier
\subsubsection*{Reglas de negocio (REGN)}

La Tabla \ref{tab:regn-inventario} agrupa las restricciones ineludibles que condicionan las operaciones del sistema, definiendo estrictamente qui\'en cuenta con los privilegios para ejecutar cada acci\'on, bajo qu\'e contexto acad\'emico particular y respetando qu\'e condiciones temporales o de visibilidad.

\begin{table}[!htbp]
\centering
\small
\begin{tabular}{|P{1.6cm}|P{3.5cm}|P{6.6cm}|}
\hline
\textbf{C\'odigo} & \textbf{Nombre} & \textbf{Regla de negocio} \\
\hline
REGN-001 & Correcci\'on de entregables & Solo el profesor asignado al entregable puede emitir nota. \\
\hline
REGN-002 & Visibilidad de entregables & El alumno solo accede a sus propias entregas/versiones. \\
\hline
REGN-003 & Correcci\'on de actividades & Solo el profesor asignado a la actividad puede calificar. \\
\hline
REGN-004 & Visibilidad de actividades & El alumno solo visualiza sus propias resoluciones de actividad. \\
\hline
REGN-005 & Fecha l\'imite de entrega & No se admiten env\'ios fuera de plazo salvo configuraci\'on expl\'icita. \\
\hline
REGN-006 & Recursos no visibles & Se deniega interacci\'on del alumno con elementos ocultos. \\
\hline
REGN-007 & Creaci\'on de recursos & Solo profesorado crea actividades y entregables. \\
\hline
REGN-008 & Emisi\'on de feedback & Solo el docente autorizado publica feedback evaluativo. \\
\hline
\end{tabular}
\normalsize
\caption{Inventario de reglas de negocio (REGN) en ERS}
\label{tab:regn-inventario}
\end{table}

\FloatBarrier
\subsubsection*{Requisitos no funcionales (RQNF)}

La Tabla \ref{tab:rqnf-inventario-completo} resume las condiciones de calidad t\'ecnica y operacionales que debe cumplir el sistema para garantizar una ejecuci\'on fiable, mantenible y segura, m\'as all\'a de la estricta funcionalidad visible por el usuario.

\begin{table}[!htbp]
\centering
\small
\begin{tabular}{|P{1.6cm}|P{3.2cm}|P{6.9cm}|}
\hline
\textbf{C\'odigo} & \textbf{Nombre} & \textbf{Condici\'on de calidad} \\
\hline
RQNF-001 & Tolerancia a fallos & Persistencia local temporal ante fallos de conectividad para evitar p\'erdida de datos. \\
\hline
RQNF-002 & Intuitividad & Interfaz comprensible para nuevos usuarios con patrones de uso habituales. \\
\hline
RQNF-003 & Uniformidad & Coherencia visual y de estilo con directrices de marca institucional. \\
\hline
RQNF-004 & Acoplamiento bajo & Separaci\'on de responsabilidades con dependencias estrictamente necesarias. \\
\hline
RQNF-005 & Cohesi\'on alta & Organizaci\'on interna de m\'odulos orientada a mantenimiento evolutivo. \\
\hline
RQNF-006 & Tiempo de carga & Tiempo de transici\'on objetivo inferior a 2 segundos en condiciones \`optimas. \\
\hline
RQNF-007 & Tiempo de b\'usqueda & Respuesta r\'apida en consultas frecuentes y minimizaci\'on de b\'usquedas redundantes. \\
\hline
RQNF-008 & Compatibilidad de navegadores & Soporte para navegadores modernos de uso com\'un. \\
\hline
RQNF-009 & Exposici\'on de informaci\'on & Visualizaci\'on de datos restringida a usuarios autenticados y autorizados. \\
\hline
RQNF-010 & Encriptaci\'on & Protecci\'on criptogr\'afica de informaci\'on personal/sensible. \\
\hline
RQNF-011 & Portabilidad web & Compatibilidad expl\'icita con Chrome, Firefox y Edge en escritorio/m\'ovil. \\
\hline
\end{tabular}
\normalsize
\caption{Inventario de requisitos no funcionales (RQNF) en ERS}
\label{tab:rqnf-inventario-completo}
\end{table}


\FloatBarrier
\section{Documento de Arquitectura de Software (DAS)}

\subsection{Fundamentaci\'on te\'orica del DAS}
El Documento de Arquitectura de Software o Documento de An\'alisis del Sistema (DAS) constituye el puente entre la especificaci\'on te\'orica y la construcci\'on real. Gen\'ericamente, su finalidad es modelar los componentes l\'ogicos, la navegabilidad, la interacci\'on entre clases y los estados de las entidades principales.

En el marco de este proyecto, el objetivo espec\'ifico del DAS es establecer el mapa visual que justifique la arquitectura cliente-servidor adoptada, las secuencias transaccionales para la entrega y evaluaci\'on, y los mockups de interfaz que gu\'ian la experiencia de usuario final.

\subsection{Reutilizaci\'on de diagramas y mockups del DAS}

Para ilustrar la arquitectura y la interacci\'on del usuario con el sistema, se incluyen a continuaci\'on los diagramas y mockups principales definidos en el DAS.

\begin{table}[!htbp]
\centering
\small
\begin{tabular}{|P{3.5cm}|P{5.5cm}|P{2.8cm}|}
\hline
\textbf{Artefacto DAS} & \textbf{Archivo original en assets} & \textbf{Uso en memoria} \\
\hline
Arquitectura l\'ogica & ArquitecturaTFG.drawio.png & Dise\~no global \\
\hline
Diagrama de clases & Diagrama de clases.png & Modelo de dominio \\
\hline
Diagrama de estado & Diagrama de estado.png & Ciclo de vida \\
\hline
Navegabilidad & Diagrama de navegabilidad (1).png & Flujo de interfaz \\
\hline
Secuencia autenticaci\'on & Autenticacion y Acceso.png & Flujo de acceso \\
\hline
Secuencia crear actividad & Backend Activity Creation-2026-03-18-161459.png & Flujo docente \\
\hline
Secuencia realizar entrega & Realizar Entrega.png & Flujo estudiante \\
\hline
Secuencia evaluaci\'on & Evaluacion y Feedback.png & Flujo de correcci\'on \\
\hline
\end{tabular}
\normalsize
\caption{Diagramas principales reutilizados desde DAS}
\label{tab:das-diagramas-principales}
\end{table}

\figura{0.90}{\detokenize{../assets/ArquitecturaTFG.drawio.png}}{Arquitectura l\'ogica del sistema (DAS)}{fig:das-arquitectura}{}

\figura{0.85}{\detokenize{../assets/Diagrama de clases.png}}{Modelo de clases del sistema (DAS)}{fig:das-clases}{}

\figura{0.55}{\detokenize{../assets/Diagrama de estado.png}}{Diagrama de estado de actividades (DAS)}{fig:das-estado}{}

\figura{0.85}{\detokenize{../assets/Diagrama de navegabilidad (1).png}}{Navegabilidad principal de la interfaz (DAS)}{fig:das-navegabilidad}{}

\figura{0.88}{\detokenize{../assets/Autenticacion y Acceso.png}}{Secuencia de autenticaci\'on y acceso (DAS)}{fig:das-seq-auth}{}

\figura{0.88}{\detokenize{../assets/Realizar Entrega.png}}{Secuencia de realizaci\'on de entrega (DAS)}{fig:das-seq-entrega}{}

\figura{0.75}{\detokenize{../assets/InterfazDeUsuarioTFG-LogueadoComoEstudiante.png}}{Mockup de panel principal de estudiante (DAS)}{fig:das-mock-estudiante}{}

\figura{0.75}{\detokenize{../assets/LogueadoComoProfesor.png}}{Mockup de panel principal de profesor (DAS)}{fig:das-mock-profesor}{}

Adicionalmente, el DAS incluye un inventario de diecis\'eis mockups que cubre
las operaciones esenciales por rol. Para mantener equilibrio entre exhaustividad
y legibilidad, su trazado completo se resume en la Tabla~\ref{tab:das-mockups}.

\begin{table}[!htbp]
\centering
\small
\begin{tabular}{|P{4.9cm}|P{6.8cm}|}
\hline
\textbf{Archivo de mockup} & \textbf{Escenario principal representado} \\
\hline
InterfazDeUsuarioTFG-InicioDeSesi\'on.drawio.png & Inicio de sesi\'on y entrada al sistema. \\
\hline
InterfazDeUsuarioTFG-Registrarse.drawio.png & Registro o alta inicial de usuario. \\
\hline
InterfazDeUsuarioTFG-LogueadoComoEstudiante.png & Panel inicial de estudiante autenticado. \\
\hline
VistaEstudianteDentroDeUna\-Asignatura.png & Vista de asignatura desde rol estudiante. \\
\hline
LogueadoComoEstudiante\-Seleccionando\-Checkbox.png & Selecci\'on de entregables por estudiante. \\
\hline
LogueadoComoEstudianteEn\-Calificaciones.png & Consulta de calificaciones y feedback. \\
\hline
LogueadoComoProfesor.png & Panel inicial de profesor autenticado. \\
\hline
VistaProfesorEnAsignatura.png & Vista de asignatura desde rol profesor. \\
\hline
Creacion de actividad.png & Creaci\'on de actividad acad\'emica. \\
\hline
Edici\'on de actividad.png & Edici\'on de actividad existente. \\
\hline
Realizar calificaciones.png & Proceso de calificaci\'on por docente. \\
\hline
DentroDeCalificacionComo\-Profesor.png & Vista detallada de correcci\'on. \\
\hline
EntregaSinNadaPuesto.png & Estado inicial de entrega sin adjuntos. \\
\hline
EntregaConAlgoPuesto.png & Estado de entrega con archivo adjunto. \\
\hline
PonerComentarios.png & Edici\'on de comentarios de evaluaci\'on. \\
\hline
PaginaEstudiantesProfesores.png & Consulta de estudiantes y trazabilidad docente. \\
\hline
\end{tabular}
\normalsize
\caption{Inventario de mockups de interfaz reutilizados desde DAS}
\label{tab:das-mockups}
\end{table}

\subsection{Lectura t\'ecnica de los diagramas estructurales del DAS}

La reutilizaci\'on de las Figuras~\ref{fig:das-arquitectura},
\ref{fig:das-clases}, \ref{fig:das-estado} y \ref{fig:das-navegabilidad}
no se limita a un uso ilustrativo, sino que establece una l\'inea de dise\~no
que condiciona decisiones de implementaci\'on en backend, frontend y datos.

En particular, la arquitectura l\'ogica permite justificar la separaci\'on entre
interacci\'on con usuario, seguridad, gesti\'on acad\'emica, procesamiento de
archivos y persistencia. Este reparto reduce acoplamiento transversal,
facilita pruebas por componente y mejora mantenibilidad evolutiva.

El modelo de clases se emplea como referencia para mantener coherencia entre
entidades persistentes, DTOs y contratos de API. Del mismo modo, el diagrama
de estado de actividades se refleja en validaciones de servicio que evitan
transiciones no permitidas (por ejemplo, publicar, ocultar o cerrar fuera de
las reglas definidas).

Por \'ultimo, el diagrama de navegabilidad refuerza la consistencia entre
estructura funcional y experiencia de usuario por rol, evitando rutas
incoherentes con las restricciones de autorizaci\'on descritas en ERS.

\begin{table}[!htbp]
\centering
\small
\begin{tabular}{|P{2.5cm}|P{3.6cm}|P{3.8cm}|P{2.4cm}|}
\hline
\textbf{Artefacto DAS} & \textbf{Decisi\'on de dise\~no derivada} & \textbf{Materializaci\'on t\'ecnica} & \textbf{Requisitos ERS reforzados} \\
\hline
Arquitectura l\'ogica & Separaci\'on funcional en subsistemas & Frontend SPA + API REST + servicios por dominio + repositorios & SUBS-001 a SUBS-005, RQNF-004 \\
\hline
Diagrama de clases & Normalizaci\'on del dominio acad\'emico & Entidades \texttt{Usuario}, \texttt{Curso}, \texttt{Actividad}, \texttt{Entregable}, \texttt{Entrega} y \texttt{Feedback} & RQI-001 a RQI-005, RQF-005 a RQF-020 \\
\hline
Diagrama de estado & Control expl\'icito del ciclo de vida & Validaciones de transici\'on y reglas temporales en servicios & REGN-005, REGN-006, RQF-006, RQF-024 \\
\hline
Navegabilidad & Flujos de interfaz por rol y contexto & Rutas protegidas, vistas diferenciadas y control de visibilidad & RQG-001, RQG-007, RQNF-002, RQNF-009 \\
\hline
\end{tabular}
\normalsize
\caption{Interpretaci\'on t\'ecnica de los diagramas estructurales del DAS}
\label{tab:das-lectura-estructural}
\end{table}

\subsection{Flujos operativos del DAS y su encaje en la soluci\'on}

Adem\'as de los diagramas estructurales, el DAS aporta una visi\'on de
comportamiento detallada mediante secuencias y mockups operativos. Esta capa
es especialmente relevante para documentar el recorrido completo de profesor y
estudiante en los procesos cr\'iticos del MVP.

El flujo de creaci\'on de actividad docente se explicita en la
Figura~\ref{fig:das-seq-crear-actividad}. En ella se observa la
coordinaci\'on entre interfaz, validaciones de negocio y persistencia,
alineada con RQF-005, RQF-008 y RQF-011.

\figura{0.88}{\detokenize{../assets/Backend Activity Creation-2026-03-18-161459.png}}{Secuencia de creaci\'on de actividad en backend (DAS)}{fig:das-seq-crear-actividad}{}

El ciclo de evaluaci\'on y retroalimentaci\'on queda recogido en la
Figura~\ref{fig:das-seq-evaluacion}, que sirve de base para justificar la
integraci\'on entre calificaci\'on, comentarios y publicaci\'on de resultados,
en coherencia con RQF-018, RQF-019 y REGN-001/REGN-008.

\figura{0.88}{\detokenize{../assets/Evaluacion y Feedback.png}}{Secuencia de evaluaci\'on y feedback (DAS)}{fig:das-seq-evaluacion}{}

En la dimensi\'on de interfaz, los mockups de trabajo muestran c\'omo se
materializa la navegaci\'on por rol en escenarios reales de uso. Las
Figuras~\ref{fig:das-mock-estudiante-asignatura} y
\ref{fig:das-mock-calificacion-profesor} capturan dos momentos clave del
producto: seguimiento de actividad por el estudiante y correcci\'on detallada
por el profesorado.

\figura{0.78}{\detokenize{../assets/VistaEstudianteDentroDeUnaAsignatura.png}}{Mockup operativo: estudiante dentro de una asignatura (DAS)}{fig:das-mock-estudiante-asignatura}{}

\figura{0.78}{\detokenize{../assets/DentroDeCalificacionComoProfesor.png}}{Mockup operativo: detalle de calificaci\'on como profesor (DAS)}{fig:das-mock-calificacion-profesor}{}

Como apoyo adicional, la Figura~\ref{fig:das-mock-entrega-con-archivo}
ilustra el estado de entrega con contenido adjunto, \`util para justificar
la trazabilidad de versiones y el control de evidencias de env\'io.

\figura{0.70}{\detokenize{../assets/EntregaConAlgoPuesto.png}}{Mockup de entrega con archivo adjunto (DAS)}{fig:das-mock-entrega-con-archivo}{}

\begin{table}[!htbp]
\centering
\small
\begin{tabular}{|P{2.1cm}|P{3.4cm}|P{4.0cm}|P{2.8cm}|}
\hline
\textbf{Actor} & \textbf{Operaci\'on central} & \textbf{Evidencia DAS reutilizada} & \textbf{Requisitos ERS asociados} \\
\hline
Administrador & Alta y administraci\'on de identidades/cursos & Mockups de alta de usuario y creaci\'on de curso & RQF-001 a RQF-004, RQF-021, RQF-022 \\
\hline
Profesor & Dise\~no y publicaci\'on de actividad/entregable & Fig. \ref{fig:das-seq-crear-actividad} y mockup de creaci\'on de actividad & RQF-005 a RQF-013, REGN-007 \\
\hline
Estudiante & Entrega y seguimiento de estado & Fig. \ref{fig:das-seq-entrega}, Fig. \ref{fig:das-mock-entrega-con-archivo} & RQF-014 a RQF-017, REGN-005 \\
\hline
Profesor evaluador & Correcci\'on y feedback & Fig. \ref{fig:das-seq-evaluacion}, Fig. \ref{fig:das-mock-calificacion-profesor} & RQF-018, RQF-019, REGN-001, REGN-008 \\
\hline
Profesor y estudiante & Consulta y trazabilidad de resultados & Mockups de comentarios y vista cruzada docente-estudiante & RQF-020, RQI-005, RQNF-009 \\
\hline
\end{tabular}
\normalsize
\caption{Cobertura operativa por rol con evidencia de DAS y requisitos de ERS}
\label{tab:das-cobertura-rol}
\end{table}

\subsection{Cat\'alogo de casos de uso y matriz de trazabilidad ERS-DAS-implementaci\'on}

El DAS define un cat\'alogo amplio de casos de uso; en esta memoria se han
seleccionado cuatro diagramas representativos (Figuras~\ref{fig:das-cu-1} a
\ref{fig:das-cu-4}) que cubren acceso, gesti\'on acad\'emica, entregas y
evaluaci\'on. Esta selecci\'on se corresponde con los flujos implementados en el
MVP y sirve como base visual para enlazar el modelo operativo con la evidencia
de implementaci\'on y pruebas.

Como representaci\'on visual de dicha cobertura, se incluyen cuatro diagramas
de casos de uso de referencia (Figuras~\ref{fig:das-cu-1} a
\ref{fig:das-cu-4}), para documentar el enlace entre el cat\'alogo funcional y
los flujos reales desarrollados.

\figura{0.82}{\detokenize{../assets/CU-1.drawio.png}}{Caso de uso representativo CU-1 (DAS)}{fig:das-cu-1}{}

\figura{0.82}{\detokenize{../assets/CU-2.drawio.png}}{Caso de uso representativo CU-2 (DAS)}{fig:das-cu-2}{}

\figura{0.82}{\detokenize{../assets/CU-3.drawio.png}}{Caso de uso representativo CU-3 (DAS)}{fig:das-cu-3}{}

\figura{0.82}{\detokenize{../assets/CU-4.drawio.png}}{Caso de uso representativo CU-4 (DAS)}{fig:das-cu-4}{}

\FloatBarrier

Para resumir la cobertura por familias operativas dentro del MVP, la
Tabla~\ref{tab:das-familias-casos-uso} agrupa el alcance funcional, el actor
dominante y el grado de cobertura logrado.

\begin{table}[!htbp]
\centering
\small
\begin{tabular}{|P{3.0cm}|P{4.0cm}|P{2.7cm}|P{2.5cm}|}
\hline
\textbf{Familia operativa} & \textbf{Alcance funcional en DAS/ERS} & \textbf{Actor dominante} & \textbf{Cobertura en MVP} \\
\hline
Acceso y sesi\'on & Inicio de sesi\'on, control de identidad, continuidad de sesi\'on & Todos los roles & Completa \\
\hline
Administraci\'on acad\'emica base & Alta/modificaci\'on de cursos, usuarios y contexto de trabajo & Administrador & Completa \\
\hline
Dise\~no docente de actividades & Creaci\'on, edici\'on, visibilidad y configuraci\'on de entregables & Profesor & Completa \\
\hline
Gesti\'on de entregas & Publicaci\'on de trabajo, versionado y consulta de historial & Estudiante & Completa \\
\hline
Evaluaci\'on y feedback & Calificaci\'on, comentarios y visibilidad de resultados & Mixto & Completa \\
\hline
Supervisi\'on y trazabilidad & Vistas cruzadas profesor-estudiante y revisi\'on del estado de entregas & Profesor & Parcial avanzada \\
\hline
\end{tabular}
\normalsize
\caption{Cobertura por familias de casos de uso del DAS dentro del alcance del MVP}
\label{tab:das-familias-casos-uso}
\end{table}

\FloatBarrier

Para reforzar la coherencia documental, la Tabla~\ref{tab:trazabilidad-ampliada}
integra tres niveles de evidencia: requisito formal en ERS, artefacto de
dise\~no en DAS y materializaci\'on en componentes software y pruebas del
proyecto.

\begin{table}[!htbp]
\centering
\small
\begin{tabular}{|P{2.5cm}|P{3.5cm}|P{3.0cm}|P{2.9cm}|}
\hline
\textbf{Bloque ERS} & \textbf{Artefacto DAS reutilizado} & \textbf{Implementaci\'on principal} & \textbf{Evidencia de verificaci\'on} \\
\hline
RQG-001, RQG-007 & Arquitectura + navegabilidad + mockups de sesi\'on & Configuraci\'on de seguridad, filtros JWT, control de rutas & Tests de autenticaci\'on/autorizaci\'on \\
\hline
RQI-001 a RQI-005 & Diagrama de clases + vistas por rol & Entidades JPA, DTOs y contratos API & Pruebas de persistencia e integraci\'on \\
\hline
RQF-005 a RQF-013 & Secuencia de creaci\'on docente + mockups de actividad & Servicios de actividad/entregable y controladores asociados & Unit tests y API tests de gesti\'on docente \\
\hline
RQF-014 a RQF-017 & Secuencia de entrega + estados de mockup & Servicio de entregas, validaciones y versionado & Integraci\'on de entregas y E2E de alumno \\
\hline
RQF-018, RQF-019 & Secuencia de evaluaci\'on + detalle de calificaci\'on & Flujo de feedback/calificaci\'on en backend y frontend & Pruebas funcionales de evaluaci\'on \\
\hline
REGN-001 a REGN-008 & Casos de uso por rol + navegabilidad & Validaciones contextuales en servicios y filtros & Pruebas de reglas de negocio \\
\hline
RQNF-002, RQNF-003 & Inventario de mockups y patrones de interfaz & Dise\~no de pantallas y mensajes de estado & E2E sobre flujos completos \\
\hline
RQNF-004, RQNF-005, RQNF-009 & Arquitectura de subsistemas + separaci\'on de capas & Estructura modular backend/frontend + control de acceso & CI/CD, test suites y auditor\'ia t\'ecnica \\
\hline
\end{tabular}
\normalsize
\caption{Trazabilidad ampliada entre ERS, DAS e implementaci\'on efectiva}
\label{tab:trazabilidad-ampliada}
\end{table}

\FloatBarrier

\subsection{Cobertura documental alcanzada}

Con la incorporaci\'on anterior, este cap\'itulo integra de forma expl\'icita los cinco bloques completos del cat\'alogo ERS (RQG, RQI, RQF, REGN y RQNF), la descomposici\'on en subsistemas usada durante dise\~no e implementaci\'on, los diagramas estructurales y de comportamiento del DAS, el cat\'alogo operativo de casos de uso y familias funcionales del sistema, el conjunto de mockups de interfaz como evidencia de dise\~no funcional y una matriz ampliada de trazabilidad entre especificaci\'on, dise\~no, implementaci\'on y verificaci\'on. Esta integraci\'on evita vac\'ios entre an\'alisis y ejecuci\'on.

De esta forma, la memoria no solo resume resultados: tambi\'en preserva la
trazabilidad documental completa entre especificaci\'on inicial, dise\~no t\'ecnico
y comportamiento implementado.

\FloatBarrier
